\subsection*{76. Education: The Product Is an Improvement in Human Capability}

Educational differentiation deserves a deeper model because a course's apparent product (lessons and curriculum) differs from its true purpose: \emph{making someone capable of accomplishing what they could not accomplish before}. Learners provide an essential input: their effort. Outcomes depend on baseline preparation, motivation, teaching, practice, feedback, social environment, and circumstances outside the institution's control.

This creates 7 particularly important sources of educational difference:

\begin{enumerate}
    \item \textbf{Starting capability:} The learner's baseline and prerequisites;
    \item \textbf{Capability destination:} The exact competence to be achieved;
    \item \textbf{Teaching and practice mechanism:} How learning is experienced and applied;
    \item \textbf{Accountability system:} Pacing, deadlines, and human review;
    \item \textbf{Credible assessment:} Verifiable demonstration of acquired judgment;
    \item \textbf{Peer environment:} Calibration against serious colleagues; and
    \item \textbf{Resource fit:} Compatibility with the learner's schedule and budget.
\end{enumerate}

For an unknown institute, course volume and superficial prestige are weak foundations. An inspectable learning process and defensible assessment can be built early. Recognized academic status, durable employer signaling, alumni influence, research reputation, and brand trust require repeated work over time. An institution cannot borrow their value from its aspirations.

\subsection*{77. Education's Uncomfortable Counterfactual}

Suppose a learner completes an advanced AI course and later receives a promotion. Did the course cause the promotion? Perhaps; or perhaps the learner was already unusually capable, the employer had planned the promotion, the learner's job performance would have improved through internal work, or the labor market shifted. An honest institute measures capability improvement without automatically claiming responsibility for every later achievement.

Differentiation should emphasize \textbf{controllable commitments and independently assessable work}: diagnostic baseline, authentic assignments, review standards, feedback timeliness, technical defenses, remediation support, and representative reporting. Learners value career upside, but marketing should separate possible future outcomes from outcomes the institution can responsibly substantiate.

\subsection*{78. The Cross-Industry Diagnostic}

When entering any unfamiliar market, analyze where customer selection friction and willingness to sacrifice originate:

\begin{table*}[t]
\centering
\small
\begin{tabularx}{\textwidth}{>{\raggedright\arraybackslash}p{5.5cm} >{\raggedright\arraybackslash}X}
\toprule
\textbf{Industry Property} & \textbf{If High or Important, Investigate...} \\
\midrule
Frequent, low-attention purchases & Recognition, mental availability, physical shelf placement, and habit \\
High-stakes, rare purchases & Verifiable proof, downside risk, committee consensus, and warranties \\
Hard-to-observe quality & Pilots, independent assessments, provenance, and trusted references \\
High switching friction & Migration tools, integration depth, and early time to first value \\
Scarce human judgment & Expert capacity, codified review procedures, and contribution per hour \\
Physical inventory and capex & Logistics, reliability, defect economics, and working-capital drag \\
Two-sided coordination & Liquidity, density, safety, matching quality, and provider earnings \\
Network and community value & Actual participant contributions and engagement, not member counts \\
Rapid technological commoditization & Workflow ownership, proprietary context, and adaptive capability \\
Institutional legitimacy & Governance, verifiable standards, professional recognition, continuity \\
\bottomrule
\end{tabularx}
\caption{Cross-industry diagnostic for evaluating customer choice mechanisms.}
\label{tab:ch8-industry-diagnostic}
\end{table*}
